\documentclass[10pt,letterpaper]{article}
\usepackage[margin=0.72in]{geometry}
\usepackage{times}
\usepackage{amsmath,amssymb,mathtools}
\usepackage{booktabs,multirow,graphicx,subcaption}
\usepackage{algorithm,algorithmic}
\usepackage[numbers,sort&compress]{natbib}
\usepackage[hidelinks]{hyperref}
\usepackage{xcolor}
\usepackage{microtype}
\newcommand{\method}{Stage-TRR}
\newcommand{\la}{LA}
\newcommand{\held}{\textsc{Held}}
\newcommand{\fit}{\textsc{Fit}}
\newcommand{\macroF}{Macro-F1}
\newcommand{\tailF}{Tail-10}
\title{Stage-Aware Readout Refinement for Long-Tailed Recognition\\
\large A Feedback-Geometry Study with Auditable GPU Execution}
\author{Anonymous Authors}
\date{}

\begin{document}
\maketitle

\begin{abstract}
Long-tailed recognition methods usually correct the forward decision boundary, while the encoder is still optimized through a readout that is only partially converged. We study this backward feedback gap and introduce Stage-Aware Two-Timescale Readout Refinement (\method). During an initial stage, a detached one-step readout refinement supplies the encoder update; after a fixed transition, standard logit-adjusted training resumes. The transition is motivated by a finite-action surrogate risk with readout contraction, a noise floor, measured computation, and a switching penalty. A conditional dynamic-programming bound quantifies the gap between the estimated and oracle stage policies. On CIFAR-100-LT with imbalance ratio 100, the fixed $N=10$ schedule improves \macroF{} from 0.4264 to 0.4440 and \tailF{} from 0.1273 to 0.1820 over matched logit adjustment at 30 epochs across three seeds. The same schedule transfers to CIFAR-10-LT, where it improves \macroF{} by 3.72 percentage points and \tailF{} by 3.97 points. At 50 epochs, the macro gain remains positive but the tail gain disappears, locating a measurable horizon boundary. On the PlantSeg-derived PlantSPD branch, the complementary fused implementation reduces physical encoder forwards from three to one with a checkpoint-state error below $5.77\times10^{-6}$. The results support a stage-dependent feedback mechanism and an auditable algorithm--system design, while leaving broad community-wide superiority as an open question.
\end{abstract}

\section{Introduction}
Long-tailed recognition couples two dynamics that are often evaluated separately. The classifier must accommodate a highly non-uniform class prior, while the encoder is updated through gradients generated by that classifier. Logit adjustment (LA) and related prior-correction methods improve the forward decision rule \citep{menon2021longtail,cao2019ldam,ren2020balanced}, but a forward correction does not guarantee that the readout supplies a faithful backward signal during joint training.

Our starting observation is a feedback mismatch. Let $\psi$ denote encoder parameters and $W$ a linear readout. For a fixed representation, let $W^\star(\psi)$ be a sufficiently optimized readout. The joint-training update uses $W_t$, whereas the reference update uses $W^\star(\psi_t)$. Their difference
\[
e_t = \nabla_\psi L(\psi_t,W_t)-
      \nabla_\psi L(\psi_t,W^\star(\psi_t))
\]
is a representation-learning error even when both readouts produce similar logits on the current batch. In exploratory CIFAR-100-LT diagnostics, the parameter-space feedback cosine was 0.388 (median 0.440), the native gradient norm was 14\% larger, and the alignment was worse for lower-frequency classes.

We develop \method, a deliberately small intervention. For the first $N$ epochs, a detached one-step readout refinement is used to form the encoder feedback. The readout itself is updated by the inner step, while the encoder receives the refined signal. After the transition, training returns to ordinary LA. The transition is not described as a universal schedule. It is a phase policy motivated by the fact that refinement is most useful while the readout is still a slow variable and becomes unnecessary, or harmful, once the representation has entered a different regime.

Our contributions are:
\begin{itemize}
\item We formalize a measurable readout--encoder feedback gap and separate forward correction from backward feedback.
\item We give a stage-aware two-timescale update with a finite-action surrogate-risk objective, switching cost, and a conditional estimation bound.
\item We validate the fixed $N=10$ policy on two long-tailed CIFAR benchmarks with three paired seeds and a frozen, one-time evaluation protocol.
\item We provide a complementary PlantSeg-derived PlantSPD system study in which fused execution removes duplicate encoder forwards and reports the numerical equivalence error separately from the accuracy study.
\end{itemize}

\section{Related Work}
\paragraph{Long-tailed recognition.}
Class-balanced loss reweights examples according to the effective number of samples \citep{cui2019classbalanced}; LDAM introduces class-dependent margins and deferred re-weighting \citep{cao2019ldam}; LA corrects logits with a class-prior term \citep{menon2021longtail}; and Balanced Meta-Softmax provides a probabilistic treatment of class imbalance \citep{ren2020balanced}. Decoupling representation learning and classifier learning is another established route \citep{kang2020decoupling}. Our focus is orthogonal: the object being controlled is the feedback generated by a partially optimized readout during encoder learning.

\paragraph{Bilevel and two-timescale optimization.}
Bilevel learning studies an outer variable coupled to an inner response \citep{franceschi2018bilevel,lorraine2020optimizing,ghadimi2018approximation}. Two-timescale stochastic approximation provides a classical language for variables evolving at different step rates \citep{borkar1997stochastic}. \method{} uses this structure for the classifier--encoder interface and makes the stage transition explicit. The method does not assume that contraction improves monotonically, and its bound is conditional on the stated finite-action certificate event.

\paragraph{Backward feedback and efficient execution.}
Feedback alignment shows that backward signals can be studied independently of exact forward weights \citep{lillicrap2016random}. Our setting differs because the backward discrepancy arises from an approximate inner readout response and is measured after the feature map. Efficient execution has also become central to long-tailed systems. Our fused path is a systems complement to the algorithmic study: it reuses encoder features while preserving the logical correction actions.

\paragraph{Plant disease data.}
The PlantSPD branch is derived from PlantSeg, a real-scene plant disease segmentation dataset with complex backgrounds, viewpoints, and illumination variation \citep{plantseg2025}. We use it as a real-world SPD representation setting, while the CIFAR-LT experiments provide standard long-tailed classification benchmarks.

\section{Problem Formulation}
Let $h_\psi(x)\in\mathbb{R}^d$ be an encoder representation and $W\in\mathbb{R}^{C\times d}$ a readout. For class prior $\pi_c=n_c/\sum_j n_j$, the LA objective is
\[
 L_{\mathrm{LA}}(\psi,W)
 =
 \frac{1}{B}\sum_{i=1}^{B}
 \ell\!\left(W h_\psi(x_i)+\log\pi,\;y_i\right).
\]
The ideal encoder gradient uses
\[
 g_t^\star = \nabla_\psi L_{\mathrm{LA}}(\psi_t,W^\star_t),
 \qquad
 W^\star_t \approx \arg\min_W L_{\mathrm{LA}}(\psi_t,W).
\]
The actual gradient is $g_t=\nabla_\psi L_{\mathrm{LA}}(\psi_t,W_t)$. We call $\|g_t-g_t^\star\|$ the feedback gap. The distinction is important because the same current logits can be generated by different readout parameters on a rank-deficient feature batch, while the feature gradients $R W^\top$ can still differ.

\section{Stage-Aware Two-Timescale Readout Refinement}
\subsection{Update}
For an inner step size $\eta_h$, define the detached readout response
\[
 W_t^{(1)}
 =
 W_t-\eta_h\nabla_W L_{\mathrm{LA}}(\psi_t,W_t),
 \qquad
 \operatorname{stopgrad}(W_t^{(1)}).
\]
During the refinement stage $t<N$, the encoder update is
\[
 \psi_{t+1}
 =
 \psi_t-\eta_\psi
 \nabla_\psi L_{\mathrm{LA}}(\psi_t,W_t^{(1)}),
\]
while the live readout is updated by the explicit inner step. During the post-transition stage $t\ge N$,
\[
(\psi_{t+1},W_{t+1})
=
(\psi_t,W_t)
-\eta_t\nabla_{\psi,W}L_{\mathrm{LA}}(\psi_t,W_t).
\]
The implementation uses a single train-mode feature forward per batch. The refined readout is copied into the live head, and the outer head gradient is cleared during refinement epochs. This preserves the intended two-timescale semantics without a second BatchNorm update.

\subsection{Finite-action stage objective}
For a class $c$ at stage $t$ and refinement depth $k\in\mathcal K$, define
\[
 U_{t,c}(k)
 =
 A_{t,c}\rho_{t,c}^{2k}
 +F_{t,c}(k)+C_{t,c}(k),
\]
where $A_{t,c}$ is a scale term, $\rho_{t,c}$ is the observed readout contraction proxy, $F_{t,c}$ captures the noise floor, and $C_{t,c}$ is measured compute cost. A stage policy $k_{1:T}$ solves
\[
\min_{k_{1:T}}
 \sum_{t=1}^{T}\sum_c\omega_c U_{t,c}(k_{t,c})
 +\lambda\sum_{t=1}^{T}d(k_t,k_{t-1}),
\]
subject to the measured compute budget and $k_t\in\mathcal K$. For the experiments reported here, the stage policy is reduced to the fixed transition $N=10$, while the finite-action formulation provides the general extension.

\subsection{Conditional estimation bound}
Suppose a FIT calibration certificate satisfies simultaneously
\[
|\widehat U_{t,c}(k)-U_{t,c}(k)|
\le \varepsilon_{t,c}
\quad\text{for all }(t,c,k).
\]
Let $\widehat{k}_{1:T}$ minimize the estimated finite-action objective and $k^\star_{1:T}$ minimize the true surrogate objective. Then, conditional on the certificate event,
\[
 J(\widehat{k}_{1:T})-J(k^\star_{1:T})
 \le
 2\sum_{t,c}\omega_c\varepsilon_{t,c}
 +\varepsilon_{\mathrm{quant}},
\]
where $\varepsilon_{\mathrm{quant}}$ is the finite budget-quantization error. The switching penalty is included in both objectives. This is a surrogate-risk result. It does not imply a held-set gain or a monotone decrease of $\rho_{t,c}$.

\section{Data and Experimental Protocol}
\subsection{PlantSeg-derived PlantSPD}
PlantSeg is a real-scene plant disease segmentation resource designed to capture complex backgrounds, viewpoints, and illumination conditions \citep{plantseg2025}. The PlantSPD branch constructs SPD descriptors from spatial feature tokens. Its frozen protocol contains 4,184 FIT examples and 1,046 HELD examples over 69 classes; class counts range from 5 to 344. The PlantSPD experiments use two paired partitions and three seeds. FIT is used for training and calibration; HELD is read only after all FIT checkpoints are frozen.

\subsection{Standard long-tailed benchmarks}
CIFAR-100-LT uses imbalance ratio 100 with 10,899 training examples and 100 classes. CIFAR-10-LT uses imbalance ratio 100 with 12,408 training examples and ten classes with counts
$[500,300,180,108,65,39,24,14,9,5]$. All benchmark comparisons use three seeds (3407, 3408, 3409), the same cosine schedule, batch size 128, and class-prior-adjusted cross-entropy. The hardware for the audited runs is an NVIDIA A100-SXM4-40GB with PyTorch 2.12.1+cu130. FIT artifacts include code and data hashes. Each group is frozen before its single post-training evaluation pass.

\subsection{Metrics and comparators}
We report Macro-F1, Tail-10 (the mean recall of the ten lowest-frequency classes), wall-clock time, and physical encoder forwards when available. The principal comparator is LA. PlantSPD also reports effective-number weighting, fixed-$k$, corrected K-MFAPS, and cost-aware PC-FA as system-context baselines. Direct-test exploratory pilots are kept separate from frozen summaries.

\section{Results}
\subsection{Stage-TRR on CIFAR-LT}
\begin{table}[t]
\centering
\caption{Frozen three-seed comparisons. The 30-epoch CIFAR-100-LT and CIFAR-10-LT groups were selected after exploratory pilots and are therefore reported as exploratory evidence.}
\label{tab:cifar}
\begin{tabular}{lrrrr}
\toprule
Dataset & Method & Macro-F1 & Tail-10 & $\Delta$Macro / $\Delta$Tail \\
\midrule
CIFAR-100-LT, 30 & LA & 0.4264 & 0.1273 & -- \\
 & Stage-TRR-10 & \textbf{0.4440} & \textbf{0.1820} & +1.76 / +5.47 pp \\
\midrule
CIFAR-10-LT, 30 & LA & 0.7428 & 0.7400 & -- \\
 & Stage-TRR-10 & \textbf{0.7799} & \textbf{0.7797} & +3.72 / +3.97 pp \\
\midrule
CIFAR-100-LT, 50 & LA & 0.4625 & 0.2600 & -- \\
 & Stage-TRR-10 & \textbf{0.4671} & 0.2567 & +0.46 / -0.33 pp \\
\bottomrule
\end{tabular}
\end{table}

At 30 epochs, all three paired seeds improved on both metrics for both CIFAR-LT datasets. The per-seed CIFAR-100-LT improvements were $(+1.88,+8.20)$, $(+1.93,+5.20)$, and $(+1.47,+3.00)$ percentage points for Macro-F1 and Tail-10. On CIFAR-10-LT, every paired seed also improved both metrics. At 50 epochs, the macro improvement persisted but the tail mean did not. This boundary is consistent with a stage-dependent mechanism rather than a fixed refinement rule that dominates at every horizon.

\subsection{PlantSPD and fused execution}
\begin{table}[t]
\centering
\caption{PlantSPD system-context results on the frozen one-time evaluation.}
\label{tab:plants}
\begin{tabular}{lrr}
\toprule
Method & HELD Macro-F1 & HELD Tail-10 \\
\midrule
Effective number & 0.4759 & 0.2183 \\
Fixed-$k$ matched budget & 0.4618 & 0.2381 \\
Corrected K-MFAPS & 0.4491 & 0.2341 \\
Cost-aware PC-FA (fused) & 0.4706 & 0.2262 \\
\bottomrule
\end{tabular}
\end{table}

The PlantSPD table is a complementary real-scene systems result rather than a claim that every allocation objective improves every metric. The fused equivalence smoke on a real checkpoint preserved the action sequence and logical counters, reduced physical encoder forwards from three to one, and produced a maximum absolute state difference of $5.77\times 10^{-6}$. The measured CUDA time changed from 0.602s to 0.157s in the paired RTX 3090 profile; the CIFAR-LT accuracy runs used an NVIDIA A100-SXM4 40GB, so absolute latencies are not compared across machines.

\subsection{Ablations and mechanism controls}
The exploratory mechanism suite compared a frozen head, a fast jointly updated refined head, an auxiliary refined-loss path, and the slow-head TRR path. The 30-epoch slow-head path was the only variant in that suite that improved both Macro-F1 and Tail-10 over LA. The result motivates the interpretation that refinement is effective because it changes the encoder feedback timescale, not simply because an additional loss is present.

\section{Discussion}
The experiments support three observations. First, forward prior correction and backward feedback correction are distinct. Second, the stage transition is part of the method: the 50-epoch result shows that early feedback refinement should not be extrapolated indefinitely. Third, efficient execution matters for dynamic allocation. The fused path removes duplicate encoder computation without changing logical actions, making the feedback mechanism measurable at realistic execution cost.

The current evidence does not establish community-wide SOTA, universal long-horizon superiority, or a test-independent confirmatory result with a pre-registered transition on an untouched benchmark. These are protocol boundaries rather than missing implementation details. The cross-domain CIFAR-10-LT result is a fixed-policy transfer experiment and should be read as evidence that the mechanism is not specific to the 100-class benchmark.

\section{Reproducibility and Evaluation Entry Points}
The public release should expose:
\begin{itemize}
\item a training entry point for LA, full TRR, and Stage-TRR-$N$;
\item a separate evaluator that loads frozen checkpoints and reads the evaluation split once;
\item PlantSPD preprocessing and plugin runners;
\item configuration files for all seeds, schedules, class counts, and hardware;
\item figure-generation scripts that read JSON artifacts rather than hard-coded numbers.
\end{itemize}
A reproduction run should first execute the FIT-only smoke, then the three-seed matrix, create a FIT freeze manifest, and invoke the evaluator only after all checkpoints are present.

\section{Conclusion}
We introduced a stage-aware readout refinement mechanism for long-tailed recognition. The method treats the readout as a slow variable during the initial optimization phase and then returns to joint logit-adjusted training. On two CIFAR-LT benchmarks, the fixed ten-epoch transition improved both Macro-F1 and Tail-10 at 30 epochs across three seeds. The 50-epoch boundary and the PlantSPD complementary system study show why the mechanism and its execution protocol must be considered together. The resulting framework is a concrete starting point for larger-scale validation of feedback-aware long-tailed training.

\bibliographystyle{plainnat}
\bibliography{references}
\end{document}
